% Edition 0.5, research synthesis and editorial truth contract.
\chapter[Theorem Frontier and Research Decisions]{The Theorem Frontier: Exact Joins, Alternative Routes, and Adversarial Decisions}
\label{ch:frontier}

\section{Why the surviving gap is not a lack of formal machinery}
The original relative SU(2) quantum system has a concrete bosonic
carrier, an irreducible original Fock representation of the forty-eight
Majoranas, sixteen source supercharges, a gauge-invariant physical
Hilbert subspace, and a native nonnegative self-adjoint Hamiltonian.
Its ninth-axis Spin stabilizer has a concrete compact connected
covering structure. The Fock \emph{operator} algebra has a source
Spin(9) automorphism preserving the actual adjoint and trace.
Separate first-moment and weighted-closure compactness theorems have
been independently strict-qualified. These are substantial inputs,
but they are not yet an eigenvalue theorem for the original
Hamiltonian.

The problem is the \emph{interface between the results}, not merely
counting unproved lemmas. A proof about an endomorphism fiber cannot
simply be substituted for a Fock-vector proof. A proof on the smooth
core does not automatically extend to the natural domain of the closed
Hamiltonian. A topological double cover is not an already transported
highest-weight decomposition of the physical $L^2$ representation.
Each such transfer requires a typed mathematical theorem.

\section{A source-exact dependency sequence}
The most informative sequence of new proofs is:
\begin{enumerate}[label=\textbf{Gate \arabic*.},leftmargin=*,itemsep=.45em]
\item \textbf{Group implementation.} Bind the literal source nine-vector
and real 48-label actions to a global continuous Fock-vector
Spin(9) unitary representation, preserving the negative central sign.
An algebraic star action on endomorphisms is an input but not the result.
\item \textbf{Gauge and charge covariance.} Prove commutation with the
actual original SU(2) Fock implementing action, not merely with its
orthogonal label matrices. Establish all sixteen source charge
covariances on the real smooth physical core, including the cubic
fermionic term, and identify the source output rotation.
\item \textbf{True reducing sectors.} Extend the same action through
the graph closure to the original physical closed form and associated
Hamiltonian. Construct physical isotypic projections that preserve
their natural domains; do not impose the projection on a different
model.
\item \textbf{Representation and field input.} Prove the genuine
Fock-vector first-weight bound (or state a distinct operator-valued
model), source-compatible Spin(8) branching, nonzero high-type physical
states, and a source-typed bound for the \emph{unnormalized} fermionic
field.
\item \textbf{Coercivity and closure.} Obtain a first-moment
inequality on a \emph{single} actual physical isotypic smooth core.
Show the form-core closure and weighted-Sobolev triple inclusion
needed by the accepted compactness theorem.
\item \textbf{Spectral consequence.} Only then prove compact
resolvent of the appropriate restricted original Hamiltonian and
establish the intended positive point spectrum. A claim about the
entire full spectrum or a mass gap would require a separate theorem.
\end{enumerate}
These are logical gates, not a promise to launch one CI workflow per
sentence. A gate can be decomposed into small reusable Lean modules
where source-level signatures permit it.

\section{The difference between vector and observable routes}
The \emph{original vector route} concerns the actual finite complex
Fock module $\mathcal F$ and hence $L^2(\Bconfiguration;\mathcal F)$.
A unitary group action on $\mathcal F$ would allow an honest
spatial-plus-fermionic action on the physical Hilbert space, subject
to its compatibility with the original gauge representation.
The \emph{observable route} works on
$\operatorname{End}_{\mathbb C}(\mathcal F)$ with the proven
star automorphism and the ordinary Hilbert--Schmidt pairing.
It avoids choosing a global Fock-vector implementation but pays for
that convenience with a different representation and a larger
first-weight bound. An identification of their \emph{spectral}
operators is not supplied by the algebra equivalence.

We should preserve both routes as mathematical tools. The
operator-valued argument may yield a useful theorem in its own right;
it must not be advertised as the original vector Hamiltonian theorem
until an explicit isometric intertwiner with compatible gauge,
charge, quadratic-form and operator domains has been constructed.

\section{A separate supersymmetry-sector route in current literature}
Two November 2026 preprints provide a useful independent comparison;
neither is part of the frozen source-qualification atlas.
Kabir Bajaj and Henry W. Lin,
\emph{Solving the SU(2) BFSS model},
\url{https://arxiv.org/abs/2611.14512} (10 November 2026,
especially Section~4.1 and Appendix~F), report numerical
positive-energy bound states and propose an existence argument
through supermultiplet selection. Alexander Michel,
Jaroslav Scheinpflug and Xi Yin,
\emph{99.9... percent of the SU(2) BFSS ground state},
\url{https://arxiv.org/abs/2611.09837} (6 November 2026),
give a variational ground-state method with rigorous bounds
derived using an \emph{auxiliary} operator whose zero-mode
space is related to the original Hamiltonian. An auxiliary
spectral gap is not a gap for the entire original Hamiltonian.

The Bajaj--Lin route must be described with its \emph{own}
representation labels. On a positive-energy supersymmetry
multiplet, the normalized charge operators are formally
\begin{equation}\label{eq:external-normalized-charges}
 C_\alpha=Q_\alpha H^{-1/2},\qquad
 \{C_\alpha,C_\beta\}=2\delta_{\alpha\beta}I.
\end{equation}
They use the $256$-dimensional positive-energy
multiplet polarization with Spin(9) decomposition
$\mathbf{44}\oplus\mathbf{84}\oplus\mathbf{128}$,
and propose a projector onto an $\mathbf{84}$ (or
$\mathbf{128}$) \emph{superspin} component. Their
analysis argues that the relevant total singlet
multiplet cannot enter the SU(2) scattering
channels under their asymptotic channel classification,
and develops test states by applying an operator to
the normalizable ground state. The preprint also
reports numerical energies for such states.

This mechanism is \emph{not} our high-$\lambda_2$
Spin(9)-isotypic oscillator exclusion.
The $\mathbf{84}$ and $\mathbf{128}$ labels refer to
components of the supersymmetry multiplet, not
a proposed splitting of the original $2^{24}$-dimensional
Fock fiber into separate generations. An overall
Spin(9) \emph{singlet} can carry these
\emph{internal superspin} labels, so no high-total-spin
coercivity inference follows from their occurrence.

\begin{openproblem}[Source-exact audit of the independent
supermultiplet argument]\label{op:external-susy-projector}
For the \emph{actual} original closed physical Hamiltonian,
verify the full source supercharge anticommutation relation
on a common appropriate domain, including the vanishing of
the Gauss term on physical vectors. Define the operators
$Q_\alpha H^{-1/2}$ first on a legitimate positive-energy
spectral subspace, verify their boundedness and Clifford
relations there, and prove that the polynomial
$\mathbf{84}$ or $\mathbf{128}$ projector is a
nonzero physical reducing projection. Independently
verify the relevant scattering-channel exclusion,
the normalizability and domain of proposed test states,
and the argument that this sector has nontrivial
unbounded positive spectrum. Finally identify all
conventions with our source Hamiltonian
$H_{\mathrm{phys}}=\tfrac1{16}T^\dagger T$.
\end{openproblem}

The appearance of $H^{-1/2}$ in
\eqref{eq:external-normalized-charges} is a domain-sensitive
operation, not a license to cancel unbounded factors
formally. On spectral subspaces separated from zero,
$H^{-1/2}$ is bounded, but the composition with the
source supercharges and the exact anticommutation
relations still require proof. A finite-dimensional
multiplet analogy alone does not establish the
closed-operator projector. In the original paper the
$\mathbf{128}$ existence modification is explicitly
left as an exercise; it is not supplied by merely
relabeling the $\mathbf{84}$ construction.

This external route deserves serious scrutiny because
it potentially changes the shortest mathematical
path to positive point spectrum. It has \emph{not}
been independently source-typed or Lean qualified in
our project, and its reported claims should not be
presented as ours. The abstract original
high-isotype compactness approach and the external
supermultiplet/channel argument remain distinct
research strategies until a precise comparison
is completed.


\section{Negative knowledge and hostile testing}
The historical FCP work established a disciplined separation between
physical counterexamples, genuinely false source equalities, Lean
elaboration failures, style/manifest errors and third-party tool
outages. For example, the large R4B/R4D source cubic assembly hit
kernel recursion limits without establishing a mathematical failure.
The NORM-0C manifest typo and deprecated consumer lemma were
infrastructure/proof-engineering errors, not counterexamples to the
operator bound. A delayed GitHub Actions job with zero workers is
not a Lean theorem failure.

Mathematical negatives deserve even greater prominence. A source
central element mapped to a \emph{scalar} operator cannot reproduce a
nontrivial sign-flip of Majorana generators by conjugation. The
operator-fiber weight interval cannot be shortened from twelve to six
by ignoring its dual input. A rank-one potential valley is not a
physical Weyl sequence. A separately normalized quadratic Gram
observable cannot automatically replace the linear physical fermionic
interaction. A color triplet is not automatically a triple of
physical particle families. These statements delimit genuine routes
that our future work should not repeat in disguise.

\section{Completion means a truthful internal edition}
An internal manuscript edition is complete when its mathematical
statements, hypotheses, provenance and open questions are accurate
at the stated source snapshot; it does not require solving the
spectral problem or a prospective unification program. The final
external monograph will be a separate editorial and scientific
product, extracted only after relevant results have been independently
verified, adversarially reviewed and placed into a coherent
mathematical theorem chain.

Our version control policy follows that distinction. The monograph
lives on an isolated documentation branch; the original Lean
sources, pinned compilation manifests, and historical receipts remain
unaltered. Local LaTeX compilation and inspection are independent of
GitHub Actions or Pages. Eventual promotion into the unified
integration branch should be a non-forced fast-forward after checking
ancestry and integrating any intervening scientific work.
